\begin{table}[H]
\centering\sffamily
\footnotesize
\setlength{\tabcolsep}{3pt}
\caption{Income concentration of delay, overall and by period of diagnosis}
\label{tab:s7}
\begin{tabular}{@{}lrrrr@{}}
\toprule
\rowcolor{tablehead}\makecell[l]{Stratum} & \makecell[r]{2010 to 2022} & \makecell[r]{2010 to 2014} & \makecell[r]{2015 to 2019} & \makecell[r]{2020 to 2022} \\
\midrule
All men & 0.062 (0.039 to 0.083) & 0.059 (0.025 to 0.092) & 0.037 (0.011 to 0.065) & 0.027 (0.003 to 0.061) \\
Low risk & 0.103 (0.068 to 0.134) & 0.092 (0.053 to 0.136) & 0.084 (0.052 to 0.121) & 0.052 (0.017 to 0.092) \\
Intermediate risk & 0.075 (0.051 to 0.096) & 0.058 (0.023 to 0.099) & 0.048 (0.019 to 0.080) & 0.039 (0.015 to 0.067) \\
High risk & 0.041 (0.017 to 0.065) & 0.028 ($-$0.011 to 0.062) & 0.020 ($-$0.004 to 0.049) & 0.015 ($-$0.013 to 0.056) \\
Unknown risk & 0.077 (0.037 to 0.128) & 0.060 (0.001 to 0.130) & 0.023 ($-$0.041 to 0.100) & 0.049 ($-$0.005 to 0.101) \\
\bottomrule
\end{tabular}
\par\smallskip
\begin{minipage}{\linewidth}\scriptsize Erreygers-corrected concentration index (95\% cluster bootstrap interval) of waiting more than 90 days, ranking men by county median household income, poorest first. Positive values mean delay is concentrated in higher-income counties. Crude: not standardised for clinical features, and it does not separate income from rurality. Indices by period were specified after an internal review (amendment 1.9).\end{minipage}
\end{table}
